\section{Results and Discussion}
\label{sec:results}

\subsection{Experimental Setup Comparison}

The present implementation evaluates the same four sensing modalities studied
in \cite{charan2022drone}: camera image, GPS position, GPS position with drone
height, and GPS position with height and distance. Accordingly, the comparison
focuses on top-$k$ beam prediction accuracy.

Table~\ref{tab:setup_comparison} summarizes the principal experimental
differences. The reference study downsampled the original 64-element
received-power vector to a 32-beam classification problem and used 8,402
training and 3,602 test samples. In contrast, the current position-family
experiments use 6,832 training, 3,416 validation, and 1,139 test samples,
while the implemented output layers contain 64 or 65 units.

\begin{table*}[!t]
\caption{Experimental Setup Comparison With the Reference Study}
\label{tab:setup_comparison}
\centering
\renewcommand{\arraystretch}{1.12}
\begin{tabular}{p{0.18\textwidth}p{0.36\textwidth}p{0.36\textwidth}}
\hline
\textbf{Aspect} & \textbf{Reference Study \cite{charan2022drone}} &
\textbf{Present Implementation} \\
\hline
Sensing modalities & Camera image, GPS position, position with height, and
position with height and distance & Same four sensing modalities \\
Beam-class definition & Original 64-element received-power vector
downsampled to 32 beam classes & Output layers contain 64 or 65 classes,
depending on the implemented model \\
Dataset split & 8,402 training and 3,602 test samples & 6,832 training,
3,416 validation, and 1,139 test samples for the position-family experiments \\
Position-model architecture & Two-hidden-layer MLP & Three hidden layers,
each containing 512 units \\
Position-model batch size & 32 & 128 \\
Missing height and distance values & Not reported as part of the reference
setup & 617 of 11,387 position-family samples imputed using training-set
medians \\
Evaluation metric & Top-$k$ beam prediction accuracy & Top-1, Top-2, Top-3,
and Top-5 beam prediction accuracy \\
\hline
\end{tabular}
\end{table*}

The present models therefore should be interpreted as a comparative
implementation rather than an exact reproduction of \cite{charan2022drone}.
In particular, the different beam-class definitions and dataset splits prevent
a strictly controlled one-to-one accuracy comparison.

\subsection{Beam Prediction Accuracy}

\begin{table}[t]
\caption{Top-$k$ Accuracy of the Present Implementation}
\label{tab:present_results}
\centering
\begin{tabular}{lcccc}
\hline
\textbf{Modality} & \textbf{Top-1} & \textbf{Top-2} &
\textbf{Top-3} & \textbf{Top-5} \\
 & \textbf{(\%)} & \textbf{(\%)} & \textbf{(\%)} & \textbf{(\%)} \\
\hline
Position & 57.59 & 82.00 & 91.92 & 97.81 \\
Position + Height & 69.71 & 89.03 & 95.17 & 99.03 \\
Position + Height + Distance & 74.19 & 89.99 & 96.05 & 98.95 \\
Image, refined ResNet-50 & \textbf{88.32} & \textbf{97.81} &
\textbf{99.56} & \textbf{99.82} \\
\hline
\end{tabular}
\end{table}

\begin{figure*}[t]
\centering
\includegraphics[width=0.94\textwidth]{figures/ieee_findings_comparison.pdf}
\caption{Beam prediction accuracy comparison. Panel (a) reports the complete
Top-$k$ results of the present implementation. Panel (b) compares only the
metrics for which the reference paper provides explicit numerical values.}
\label{fig:findings_comparison}
\end{figure*}

As shown in Table~\ref{tab:present_results}, the image-based model achieves
the strongest Top-1 performance, followed by the position-height-distance and
position-height models. The position-only model produces the lowest Top-1
accuracy. This trend is also illustrated in
Fig.~\ref{fig:findings_comparison}. The ranking is consistent with
\cite{charan2022drone}, which
reported that vision-aided beam prediction outperformed the position-based
alternatives and that position alone was insufficient for reliable drone beam
prediction.

\begin{table}[t]
\caption{Direct Comparison With Numeric Results Reported in
\cite{charan2022drone}}
\label{tab:reference_comparison}
\centering
\begin{tabular}{lccc}
\hline
\textbf{Metric} & \textbf{Reference} & \textbf{Present} &
\textbf{Difference} \\
 & \textbf{(\%)} & \textbf{(\%)} & \textbf{(p.p.)} \\
\hline
Image Top-1 & 86.32 & \textbf{88.32} & \textbf{+2.00} \\
Image Top-3 & 99.41 & \textbf{99.56} & \textbf{+0.15} \\
Image Top-5 & 99.69 & \textbf{99.82} & \textbf{+0.13} \\
Position Top-1 & $\approx$59 & 57.59 & $\approx$-1.41 \\
\hline
\end{tabular}
\end{table}

The refined image model exceeds the numerical image results reported in
\cite{charan2022drone} for all directly comparable metrics. Its Top-1
accuracy improves from 86.32\% to 88.32\%, corresponding to an absolute gain
of 2.00 percentage points. The Top-3 and Top-5 gains are smaller because both
implementations already operate near saturation at these values. The
position-only result of 57.59\% is close to, but below, the approximately 59\%
result reported in \cite{charan2022drone}. This agreement supports the
reference study's conclusion that practical GPS position alone does not fully
characterize the optimal beam for a drone with three-dimensional mobility.

\subsection{Effect of Height and Distance}

\begin{table}[t]
\caption{Top-1 Improvement Relative to Position-Only Prediction}
\label{tab:modality_gain}
\centering
\begin{tabular}{lcc}
\hline
\textbf{Added Modality} & \textbf{Absolute Gain} & \textbf{Relative Gain} \\
 & \textbf{(p.p.)} & \textbf{(\%)} \\
\hline
Height & +12.12 & +21.04 \\
Height and Distance & +16.59 & +28.81 \\
\hline
\end{tabular}
\end{table}

Adding height increases Top-1 accuracy from 57.59\% to 69.71\%, yielding a
12.12-percentage-point improvement. This result lies within the approximate
10--14-percentage-point gain for combined modalities reported in
\cite{charan2022drone}. Adding distance further increases Top-1 accuracy to
74.19\%, producing a 16.59-percentage-point gain over position alone and a
4.48-percentage-point gain over position-height.

\begin{figure}[t]
\centering
\includegraphics[width=\columnwidth]{figures/top1_modality_comparison.pdf}
\caption{Top-1 accuracy improvement as height, distance, and image sensing are
introduced.}
\label{fig:top1_modality_comparison}
\end{figure}

The improvement demonstrates that height and distance provide complementary
geometric information that is not sufficiently represented by noisy
two-dimensional GPS coordinates. Compared with position-height, adding
distance improves Top-1, Top-2, and Top-3 accuracy by 4.48, 0.97, and 0.88
percentage points, respectively. Its Top-5 accuracy decreases by only 0.09
percentage points, corresponding approximately to one sample in the current
test set; this difference should not be interpreted as statistically
significant.

\subsection{Vision-Aided Prediction}

The refined ResNet-50 achieves 88.32\% Top-1 accuracy, exceeding the strongest
position-family model by 14.14 percentage points and the position-only model
by 30.73 percentage points. This result reinforces the principal finding of
\cite{charan2022drone}: an RGB image captures the drone's apparent location
and orientation within the base-station field of view more effectively than
individual scalar sensing measurements. The image model also achieves 99.56\%
Top-3 and 99.82\% Top-5 accuracy. Consequently, restricting beam training to
the model's five highest-ranked candidates would retain the correct beam for
nearly all tested samples.

The present image model was refined from its best checkpoint using selective
unfreezing, frozen batch-normalization statistics, shuffled training samples,
and checkpoint averaging. This refinement increased test Top-1 accuracy from
87.62\% to 88.32\%, an absolute gain of 0.70 percentage points, while
preserving the previous checkpoint whenever validation performance did not
improve.

\begin{figure}[t]
\centering
\includegraphics[width=\columnwidth]{figures/image_model_refinement.pdf}
\caption{Image-model refinement and direct comparison with the image metrics
reported in \cite{charan2022drone}.}
\label{fig:image_model_refinement}
\end{figure}

\subsection{Principal Findings and Limitations}

The comparative evaluation supports the following findings. First,
position-only prediction remains the weakest sensing-aided approach,
confirming that two-dimensional GPS data alone is inadequate for robust drone
beam selection. Second, height and distance materially improve Top-1
accuracy, validating the conclusion in \cite{charan2022drone} that additional
geometric sensing information is important for three-dimensional drone
mobility. Third, vision-aided prediction remains the strongest modality and
exceeds the reference image Top-1 accuracy by 2.00 percentage points in the
present implementation. Finally, Top-3 and Top-5 image accuracies are close to
100\% in both studies, indicating that sensing-aided prediction can
substantially reduce exhaustive beam-training overhead.

The observed differences should not be presented as evidence of statistical
superiority over \cite{charan2022drone} without additional controlled
experiments. The implementations differ in beam-class definition, dataset
split, model depth, batch size, image initialization strategy, optimization
procedure, and treatment of missing sensor values. In addition, the present
results are based on saved single-run outputs; repeated trials with fixed
splits and confidence intervals have not yet been reported. Therefore, the
defensible conclusion is that the present findings reproduce the qualitative
modality ranking reported in \cite{charan2022drone} and achieve strong
quantitative results under a modified experimental configuration.

% Add this entry to the paper's bibliography.
\begin{thebibliography}{1}
\bibitem{charan2022drone}
G. Charan, A. Hredzak, C. Stoddard, B. Berrey, M. Seth, H. Nunez, and
A. Alkhateeb, ``Towards real-world 6G drone communication: Position and
camera aided beam prediction,'' in \emph{Proc. IEEE Global Communications
Conf. (GLOBECOM)}, 2022, pp. 2951--2956,
doi: 10.1109/GLOBECOM48099.2022.10000718.
\end{thebibliography}
